\section{Resultados}\label{sec:resultados}

\begin{notabox}{Estado desta seção}
Esta versão do relatório \textbf{não contém resultados de modelos}. As tabelas abaixo fixam o formato de relato;
as células marcadas com \lacuna\ serão preenchidas após os experimentos --- nenhum valor foi estimado ou
inventado. As estatísticas já obtidas, que são reais, estão nas Tabelas~\ref{tab:grupos}, \ref{tab:exemplos},
\ref{tab:deriva}, \ref{tab:codificacao}, \ref{tab:elegib} e~\ref{tab:blocos} e no Apêndice~\ref{apx:maquinas}.
\end{notabox}

\subsection{E1: capacidade e recuperação}\label{sec:res-e1}
A Tabela~\ref{tab:res-e1a} reúne a validação da implementação em padrões sintéticos (E1a). O experimento E1b
usa o mesmo formato com os protótipos do SMD, para $K\in\{2,5,10,20,30\}$ e $b\in\{2,4,8,16\}$. Esperamos que
a fração de pontos fixos caia perto de $K/N\approx0{,}14$ para a regra de Hebb (Tabela~\ref{tab:capacidade}).

\begin{table}[htbp]
\centering
\caption{E1a --- fração de padrões aleatórios, não viesados, que são pontos fixos (\%), por $N$, razão $K/N$ e
regra de armazenamento (H = Hebb; S = Storkey). A preencher.}
\label{tab:res-e1a}
\footnotesize
\setlength{\tabcolsep}{4pt}
\begin{tabular}{rcccccccc}
\toprule
 & \multicolumn{2}{c}{$N=67$} & \multicolumn{2}{c}{$N=122$} & \multicolumn{2}{c}{$N=224$} & \multicolumn{2}{c}{$N=412$}\\
\cmidrule(lr){2-3}\cmidrule(lr){4-5}\cmidrule(lr){6-7}\cmidrule(lr){8-9}
$K/N$ & H & S & H & S & H & S & H & S\\
\midrule
0,05  & \lacuna & \lacuna & \lacuna & \lacuna & \lacuna & \lacuna & \lacuna & \lacuna\\
0,10  & \lacuna & \lacuna & \lacuna & \lacuna & \lacuna & \lacuna & \lacuna & \lacuna\\
0,138 & \lacuna & \lacuna & \lacuna & \lacuna & \lacuna & \lacuna & \lacuna & \lacuna\\
0,20  & \lacuna & \lacuna & \lacuna & \lacuna & \lacuna & \lacuna & \lacuna & \lacuna\\
0,30  & \lacuna & \lacuna & \lacuna & \lacuna & \lacuna & \lacuna & \lacuna & \lacuna\\
\bottomrule
\end{tabular}
\end{table}

\subsection{E2: detecção de anomalias}\label{sec:res-e2}
\paragraph{Protocolo S.} A Tabela~\ref{tab:res-s} apresenta as métricas médias entre as 18 máquinas da
classe S em $L_{\mathrm{te}}$ (macro; a AUC-ROC e a PR-AUC são calculadas por máquina e promediadas), e a
Tabela~\ref{tab:res-conf} as matrizes de confusão agregadas (micro). Os modelos supervisionados são uma
referência superior (Seção~\ref{sec:baselines}).

\begin{table}[htbp]
\centering
\caption{Protocolo S --- métricas médias entre máquinas em $L_{\mathrm{te}}$ (macro). Limiar escolhido por $F_1$
em $L_{\mathrm{val}}$ agregado. Tempo: inferência por 100 amostras. A preencher.}
\label{tab:res-s}
\footnotesize
\setlength{\tabcolsep}{3pt}
\begin{tabularx}{\textwidth}{@{}Lccccccccc@{}}
\toprule
Modelo & Acur. & $P$ & $R$ & $F_1$ & $F_2$ & AUC-ROC & PR-AUC & MCC & s/100\\
\midrule
Hopfield clássico (Hebb)   & \lacuna & \lacuna & \lacuna & \lacuna & \lacuna & \lacuna & \lacuna & \lacuna & \lacuna\\
Hopfield clássico (Storkey)& \lacuna & \lacuna & \lacuna & \lacuna & \lacuna & \lacuna & \lacuna & \lacuna & \lacuna\\
Hopfield moderno           & \lacuna & \lacuna & \lacuna & \lacuna & \lacuna & \lacuna & \lacuna & \lacuna & \lacuna\\
OC-SVM                     & \lacuna & \lacuna & \lacuna & \lacuna & \lacuna & \lacuna & \lacuna & \lacuna & \lacuna\\
\emph{k}-NN (dist. ao $k$-ésimo vizinho) & \lacuna & \lacuna & \lacuna & \lacuna & \lacuna & \lacuna & \lacuna & \lacuna & \lacuna\\
\emph{Random Forest} (supervisionado) & \lacuna & \lacuna & \lacuna & \lacuna & \lacuna & \lacuna & \lacuna & \lacuna & \lacuna\\
SVM-RBF (supervisionado)   & \lacuna & \lacuna & \lacuna & \lacuna & \lacuna & \lacuna & \lacuna & \lacuna & \lacuna\\
\bottomrule
\end{tabularx}
\end{table}

\begin{table}[htbp]
\centering
\caption{Protocolo S --- matrizes de confusão agregadas (micro) em $L_{\mathrm{te}}$ (positivo = anômalo). A
preencher.}
\label{tab:res-conf}
\small
\begin{tabular}{lrrrr}
\toprule
Modelo & $TP$ & $FP$ & $FN$ & $TN$\\
\midrule
Hopfield clássico (melhor configuração) & \lacuna & \lacuna & \lacuna & \lacuna\\
Hopfield moderno                        & \lacuna & \lacuna & \lacuna & \lacuna\\
OC-SVM                                  & \lacuna & \lacuna & \lacuna & \lacuna\\
\emph{k}-NN                             & \lacuna & \lacuna & \lacuna & \lacuna\\
\emph{Random Forest}                    & \lacuna & \lacuna & \lacuna & \lacuna\\
SVM-RBF                                 & \lacuna & \lacuna & \lacuna & \lacuna\\
\bottomrule
\end{tabular}
\end{table}

\paragraph{Protocolo U.} A Tabela~\ref{tab:res-u} traz os modelos \emph{one-class} avaliados nas 28 máquinas,
com limiar sem rótulos definido por $\alpha$. A AUC-ROC e a PR-AUC independem de $\alpha$.

\begin{table}[htbp]
\centering
\caption{Protocolo U --- 28 máquinas, toda a metade rotulada, limiar sem rótulos (quantil $1-\alpha$ em
$S_{\mathrm{cal}}$). $F_1^{\ast}$: $F_1$ ótimo com acesso aos rótulos (teto otimista, não comparável). A
preencher.}
\label{tab:res-u}
\footnotesize
\setlength{\tabcolsep}{4pt}
\begin{tabular}{llccccccc}
\toprule
Modelo & $\alpha$ & Acur. & $P$ & $R$ & $F_1$ & AUC-ROC & PR-AUC & $F_1^{\ast}$\\
\midrule
\multirow{4}{*}{Hopfield clássico} & 0,5\,\% & \lacuna & \lacuna & \lacuna & \lacuna & \multirow{4}{*}{\lacuna} & \multirow{4}{*}{\lacuna} & \multirow{4}{*}{\lacuna}\\
 & 1\,\% & \lacuna & \lacuna & \lacuna & \lacuna & & & \\
 & 2\,\% & \lacuna & \lacuna & \lacuna & \lacuna & & & \\
 & 5\,\% & \lacuna & \lacuna & \lacuna & \lacuna & & & \\
\midrule
\multirow{4}{*}{Hopfield moderno} & 0,5\,\% & \lacuna & \lacuna & \lacuna & \lacuna & \multirow{4}{*}{\lacuna} & \multirow{4}{*}{\lacuna} & \multirow{4}{*}{\lacuna}\\
 & 1\,\% & \lacuna & \lacuna & \lacuna & \lacuna & & & \\
 & 2\,\% & \lacuna & \lacuna & \lacuna & \lacuna & & & \\
 & 5\,\% & \lacuna & \lacuna & \lacuna & \lacuna & & & \\
\midrule
\multirow{4}{*}{OC-SVM} & 0,5\,\% & \lacuna & \lacuna & \lacuna & \lacuna & \multirow{4}{*}{\lacuna} & \multirow{4}{*}{\lacuna} & \multirow{4}{*}{\lacuna}\\
 & 1\,\% & \lacuna & \lacuna & \lacuna & \lacuna & & & \\
 & 2\,\% & \lacuna & \lacuna & \lacuna & \lacuna & & & \\
 & 5\,\% & \lacuna & \lacuna & \lacuna & \lacuna & & & \\
\midrule
\multirow{4}{*}{\emph{k}-NN} & 0,5\,\% & \lacuna & \lacuna & \lacuna & \lacuna & \multirow{4}{*}{\lacuna} & \multirow{4}{*}{\lacuna} & \multirow{4}{*}{\lacuna}\\
 & 1\,\% & \lacuna & \lacuna & \lacuna & \lacuna & & & \\
 & 2\,\% & \lacuna & \lacuna & \lacuna & \lacuna & & & \\
 & 5\,\% & \lacuna & \lacuna & \lacuna & \lacuna & & & \\
\bottomrule
\end{tabular}
\end{table}

\subsection{Ablações do detector de Hopfield}\label{sec:res-abl}
A Tabela~\ref{tab:res-abl} mostra o efeito do número de níveis $b$, do número de protótipos $K$ e da regra de
armazenamento sobre a PR-AUC média (Protocolo~S, escore $S_1$), e a Tabela~\ref{tab:res-escores} compara os
escores da Tabela~\ref{tab:escores}.

\begin{table}[htbp]
\centering
\caption{Ablação --- PR-AUC média (macro) do detector clássico, escore $S_1$, Protocolo~S, por $b$, $K$ e regra.
A preencher.}
\label{tab:res-abl}
\small
\setlength{\tabcolsep}{5pt}
\begin{tabular}{llccccc}
\toprule
Regra & $b$ & $K=2$ & $K=5$ & $K=10$ & $K=20$ & $K=30$\\
\midrule
\multirow{4}{*}{Hebb}    & 2  & \lacuna & \lacuna & \lacuna & \lacuna & \lacuna\\
                         & 4  & \lacuna & \lacuna & \lacuna & \lacuna & \lacuna\\
                         & 8  & \lacuna & \lacuna & \lacuna & \lacuna & \lacuna\\
                         & 16 & \lacuna & \lacuna & \lacuna & \lacuna & \lacuna\\
\midrule
\multirow{4}{*}{Storkey} & 2  & \lacuna & \lacuna & \lacuna & \lacuna & \lacuna\\
                         & 4  & \lacuna & \lacuna & \lacuna & \lacuna & \lacuna\\
                         & 8  & \lacuna & \lacuna & \lacuna & \lacuna & \lacuna\\
                         & 16 & \lacuna & \lacuna & \lacuna & \lacuna & \lacuna\\
\bottomrule
\end{tabular}
\end{table}

\begin{table}[htbp]
\centering
\caption{Comparação dos escores $S_1$--$S_6$ na melhor configuração do detector clássico (Protocolo~S, macro).
A preencher.}
\label{tab:res-escores}
\small
\begin{tabular}{lcc}
\toprule
Escore & AUC-ROC & PR-AUC\\
\midrule
$S_1$ & \lacuna & \lacuna\\
$S_2$ & \lacuna & \lacuna\\
$S_3$ & \lacuna & \lacuna\\
$S_4$ & \lacuna & \lacuna\\
$S_5$ & \lacuna & \lacuna\\
$S_6$ & \lacuna & \lacuna\\
\bottomrule
\end{tabular}
\end{table}

\subsection{E3: Etapa II --- LLM \emph{versus} modelos clássicos}\label{sec:res-e3}
A Tabela~\ref{tab:res-llm} segue o protocolo da Seção~\ref{sec:etapa2}. Os modelos clássicos são reavaliados na
\emph{mesma} subamostra.

\begin{table}[htbp]
\centering
\caption{Etapa II --- LLM (modelo, versão e data: \apreencher) e modelos clássicos na mesma subamostra de
$L_{\mathrm{te}}$ (máquinas 1-3, 2-5 e 3-6; $n=\ $\apreencher). Precisão ajustada por prevalência. Custo:
estimativa de API por 100 amostras (preços \apreencher). A preencher.}
\label{tab:res-llm}
\footnotesize
\setlength{\tabcolsep}{3pt}
\begin{tabularx}{\textwidth}{@{}Lcccccccc@{}}
\toprule
Configuração & Acur. & $P_{\mathrm{nat}}$ & $R$ & $F_1$ & $F_2$ & AUC-ROC & s/100 & Custo/100\\
\midrule
LLM, \emph{zero-shot}       & \lacuna & \lacuna & \lacuna & \lacuna & \lacuna & \lacuna & \lacuna & \lacuna\\
LLM, \emph{few-shot} $k=3$  & \lacuna & \lacuna & \lacuna & \lacuna & \lacuna & \lacuna & \lacuna & \lacuna\\
LLM, \emph{few-shot} $k=5$  & \lacuna & \lacuna & \lacuna & \lacuna & \lacuna & \lacuna & \lacuna & \lacuna\\
LLM, \emph{few-shot} $k=10$ & \lacuna & \lacuna & \lacuna & \lacuna & \lacuna & \lacuna & \lacuna & \lacuna\\
LLM, \emph{few-shot} $k=20$ & \lacuna & \lacuna & \lacuna & \lacuna & \lacuna & \lacuna & \lacuna & \lacuna\\
\midrule
\emph{Random Forest}        & \lacuna & \lacuna & \lacuna & \lacuna & \lacuna & \lacuna & \lacuna & ---\\
SVM-RBF                     & \lacuna & \lacuna & \lacuna & \lacuna & \lacuna & \lacuna & \lacuna & ---\\
Hopfield clássico           & \lacuna & \lacuna & \lacuna & \lacuna & \lacuna & \lacuna & \lacuna & ---\\
OC-SVM                      & \lacuna & \lacuna & \lacuna & \lacuna & \lacuna & \lacuna & \lacuna & ---\\
\bottomrule
\end{tabularx}
\end{table}

\subsection{E4 (opcional): interpretabilidade}
Quando realizado, comparará os atributos cujos bits foram alterados na recuperação com os atributos de
\texttt{interpretation\_label}: \apreencher.
